% =========================================================================
% Theoretical Analysis of Active Spiking Perception (ASP) and the
% Slice Selection Policy (SSP) -- paper-ready section.
%
% Standalone-compilable; to insert into the AAAI template, copy everything
% between \begin{document} and \end{document} and keep the theorem
% environments (or map them onto the template's).
%
% Every result carries an explicit "Empirical verification" protocol naming
% the script (verification/V*.py) and experiment config (configs/*) that
% tests its prediction. Assumptions are stated, not hidden.
% =========================================================================
\documentclass[11pt]{article}
\usepackage{amsmath,amssymb,amsthm,mathtools}
\usepackage[margin=1in]{geometry}
\usepackage{booktabs}
\newtheorem{theorem}{Theorem}
\newtheorem{proposition}{Proposition}
\newtheorem{corollary}{Corollary}
\newtheorem{assumption}{Assumption}
\newtheorem{remark}{Remark}
\newcommand{\E}{\mathbb{E}}
\newcommand{\Prob}{\mathbb{P}}
\newcommand{\R}{\mathbb{R}}
\newcommand{\ssp}{\mathrm{ssp}}

\begin{document}

\section{Theoretical Analysis}
\label{sec:theory}

\subsection{Setup and Notation}
An input is partitioned into $K$ spatial regions (slices) with precomputed
geometry descriptors $g_m\in\R^6$, $m\in[K]$. At step $t$ the network holds
LIF membrane state $u_{t-1}\in\R^D$, the SSP scores unvisited anchors
\begin{equation}
  \mathrm{score}_t(m)\;=\;\frac{(W_k u_{t-1})^{\!\top}(W_q g_m)}{\sqrt{d_{\ssp}}},
  \qquad W_k\in\R^{d_{\ssp}\times D},\; W_q\in\R^{d_{\ssp}\times 6},
  \label{eq:score}
\end{equation}
visited anchors are masked to $-\infty$, the argmax slice is encoded into
$e_{m^\ast_t}\in\R^D$, and the head emits logits $y_t$ with softmax
$p_t\in\Delta^{C-1}$. Let the \emph{margin process} be
$M_t = p_t^{(1)}-p_t^{(2)}$ (top-1 minus top-2 probability), and define the
exit time $T_\theta=\min\{t: M_t>\theta\}\wedge K$. Let
$\mathcal{F}_t$ be the filtration generated by the first $t$ observations,
$h_t$ the full observation history, and $y$ the label.

% -----------------------------------------------------------------------
\subsection{Why margin exiting is safe: a selective-risk guarantee}

\begin{assumption}[Margin calibration]\label{as:calib}
For exit events, the top-1 softmax probability is (weakly) faithful:
$\Prob(\hat y_t = y \mid \mathcal{F}_t,\, M_t>\theta) \ge p_t^{(1)}$
in expectation over exited samples. We \emph{audit} this assumption rather
than take it on faith: expected calibration error (ECE), reliability
diagrams, and temperature-scaled variants are logged by
\texttt{S5\_calibration}.
\end{assumption}

\begin{theorem}[Selective risk of margin exiting]\label{thm:risk}
For a $C$-class problem, on the event $\{M_t \ge \theta\}$ the top-1
probability satisfies $p_t^{(1)} \ge \frac{1+(C-1)\theta}{C}$. Consequently,
under Assumption~\ref{as:calib}, the error rate among samples that exit at
margin threshold $\theta$ obeys
\begin{equation}
  \mathrm{Risk}(\theta)\;\le\;\frac{(C-1)(1-\theta)}{C}.
  \label{eq:riskbound}
\end{equation}
\end{theorem}
\begin{proof}
The remaining mass $1-p^{(1)}$ is spread over $C-1$ classes, so its maximum
is at least its mean: $p^{(2)} \ge \frac{1-p^{(1)}}{C-1}$. Combining with
$p^{(1)} \ge p^{(2)}+\theta$ gives
$p^{(1)} \ge \frac{1-p^{(1)}}{C-1}+\theta$, i.e.\
$p^{(1)}\frac{C}{C-1} \ge \frac{1}{C-1}+\theta$, hence
$p^{(1)} \ge \frac{1+(C-1)\theta}{C}$. Under Assumption~\ref{as:calib} the
conditional error is at most $1-p^{(1)} \le \frac{(C-1)(1-\theta)}{C}$.
\end{proof}

\begin{remark}
Bound~\eqref{eq:riskbound} is worst-case over confidence profiles; empirical
selective risk is typically far lower, and the gap is exactly the
calibration slack that S5 measures. The bound also explains why a
\emph{single} threshold transfers across datasets of different $C$:
at $\theta=0.7$ it guarantees $\le 29.3\%$ conditional error for $C=40$ and
$\le 27.0\%$ for $C=10$ \emph{before} calibration tightening.
\end{remark}

\noindent\textbf{Empirical verification.} \texttt{verification/V2\_selective\_risk.py}
checks \eqref{eq:riskbound} at every $\theta$ on the full grid and the
monotonicity of $\mathrm{Risk}(\theta)$; \texttt{A1\_theta} logs
\texttt{risk\_exited} against \texttt{risk\_bound\_thm1} on every benchmark
(ModelNet10/40, ScanObjectNN, CIFAR-10/100 patch-slices).

% -----------------------------------------------------------------------
\subsection{Why 2.4 slices suffice: a stopping-time law}

\begin{assumption}[Informative-drift]\label{as:drift}
Under the trained SSP, the margin process is a submartingale with drift
$\delta>0$ until exit: $\E[M_t-M_{t-1}\mid\mathcal{F}_{t-1}]\ge\delta$ for
$t< T_\theta$, with one-step increments bounded by $c\le 1$.
This is a \emph{measured} property of the trained model (drift statistics
are logged), not an a priori axiom.
\end{assumption}

\begin{theorem}[Expected slices to exit]\label{thm:stop}
Under Assumption~\ref{as:drift},
\begin{equation}
  \E[T_\theta]\;\le\;1+\frac{\big(\theta+c-\E[M_1]\big)^{+}}{\delta},
  \label{eq:wald}
\end{equation}
and $\E[T_\theta]$ is (to first order) affine in $\theta$ with slope
$1/\delta$.
\end{theorem}
\begin{proof}
$Z_t = M_t-\delta (t-1)$ is a submartingale by Assumption~\ref{as:drift};
$T_\theta$ is a bounded stopping time, so optional stopping gives
$\E[M_{T_\theta}]-\delta(\E[T_\theta]-1)\ge \E[M_1]$. On exit paths the
overshoot is at most $c$ ($M_{T_\theta}\le\theta+c$); on censored paths
$M_{T_\theta}\le\theta$. Hence
$\E[T_\theta]\le 1+(\theta+c-\E[M_1])/\delta$. Truncation at $K$ only
decreases the left side.
\end{proof}

\begin{corollary}[The $\theta=0.7$ operating point]\label{cor:24}
The reported exit at $2.4/16$ slices for $\theta=0.7$ is consistent with
\eqref{eq:wald} iff the measured drift satisfies
$\delta \gtrsim (0.7+c-\E[M_1])/1.4$. This turns the headline efficiency
number into a \emph{falsifiable} statement about two measurable scalars
$(\delta, \E[M_1])$ rather than an unexplained empirical accident.
\end{corollary}

\begin{corollary}[Failure mode: censoring and bimodality]\label{cor:bimodal}
Samples whose conditional drift is non-positive (ambiguous inputs; corrupted
regions) exit at $T_\theta=K$ with high probability. The exit-time histogram
is therefore predicted to be \emph{bimodal} --- a mass near
$1+(\theta+c-\E[M_1])/\delta$ and a censored spike at $K$ --- with the
censored mass upper-bounded by the fraction of non-positive-drift samples.
\end{corollary}

\noindent\textbf{Empirical verification.} \texttt{V1\_stopping\_time.py}
estimates $\hat\delta$, $\widehat{\E[M_1]}$, verifies
$\E[T_\theta]\le$ \eqref{eq:wald} across the $\theta$ grid, tests
affinity-in-$\theta$, and checks censored mass against negative-drift mass;
\texttt{S6\_exit\_audit} tests bimodality and per-class exit/error coupling.

% -----------------------------------------------------------------------
\subsection{Why greedy slice selection is near-optimal}

Let $\phi(S)$ denote features of a slice subset $S\subseteq[K]$ and
$V(S)=I\big(y;\phi(S)\big)$ the aggregated information gain.

\begin{assumption}[Conditional independence surrogate]\label{as:ci}
Slice features are conditionally independent given the label,
$\phi(i)\perp\phi(j)\mid y$. This naive-Bayes-style surrogate is standard
for tractable analysis of information maximization and is
\emph{empirically audited} (violation rate of the submodular inequality is
logged; see below).
\end{assumption}

\begin{theorem}[Monotone submodular information gain]\label{thm:submod}
$V$ is monotone: $V(S)\le V(S\cup\{s\})$ always. Under
Assumption~\ref{as:ci}, $V$ is submodular:
$V(S\cup\{s\})-V(S)\;\ge\;V(T\cup\{s\})-V(T)$ for all $S\subseteq T$,
$s\notin T$.
\end{theorem}
\begin{proof}
Monotonicity: $I(y;\phi(S\cup\{s\}))-I(y;\phi(S)) =
I(y;\phi(s)\mid\phi(S))\ge 0$. Submodularity: with conditional
independence given $y$,
$I(y;\phi(s)\mid \phi(S)) = H(\phi(s)\mid\phi(S)) - H(\phi(s)\mid y)$
(the second term loses its conditioning on $\phi(S)$ by
Assumption~\ref{as:ci}). Conditioning reduces entropy, so
$H(\phi(s)\mid\phi(S))\ge H(\phi(s)\mid\phi(T))$ for $S\subseteq T$, while
$H(\phi(s)\mid y)$ is common to both sides; the marginal gain is therefore
non-increasing in the conditioning set (Krause \& Guestrin, 2005).
\end{proof}

\begin{corollary}[Greedy near-optimality]\label{cor:greedy}
The greedy policy that observes
$\arg\max_s V(S\cup\{s\})-V(S)$ achieves
$V(S^{\mathrm{greedy}}_k)\ge(1-e^{-k/k^\ast})\,V(S^\ast_{k^\ast})$ against
any budget-$k^\ast$ optimum (Nemhauser et al., 1978). The SSP is an
\emph{amortized} greedy policy: Eq.~\eqref{eq:score} is trained end-to-end
so that its ranking tracks the marginal-gain ranking without computing $V$.
\end{corollary}

\begin{remark}[Honesty about Assumption~\ref{as:ci}]
Mutual information is \emph{not} submodular in general. We therefore report
the measured violation rate of the submodular inequality on sampled chains;
the greedy guarantee should be read as ``holds up to the measured violation
mass,'' which is the strongest defensible claim at review time.
\end{remark}

\noindent\textbf{Empirical verification.} \texttt{V3\_submodularity.py}:
(a) diminishing greedy marginal gains, (b) sampled-chain violation rate,
(c) anytime-accuracy ordering oracle-greedy $\ge$ SSP $\ge$ random with the
SSP--greedy gap $\ll$ random--greedy gap (amortization quality).
Full-scale protocol: \texttt{A5\_policy} anytime curves on all benchmarks.

% -----------------------------------------------------------------------
\subsection{Why the membrane state is enough input for the policy}

\begin{proposition}[Approximate sufficiency of $u_t$]\label{prop:suff}
$u_t$ is a deterministic causal filter of the history $h_t$
($y\to h_t\to u_t$ is a Markov chain), hence
$I(y;u_t)\le I(y;h_t)$ (data processing). Define the sufficiency gap
$\varepsilon_t = I(y;h_t)-I(y;u_t)\ge 0$ and beliefs
$b_t=\Prob(y\mid h_t)$, $\hat b_t=\Prob(y\mid u_t)$; then
$\E\,\mathrm{KL}(b_t\,\|\,\hat b_t)=\varepsilon_t$, and for any bounded
policy-value functional $V^\pi$ that is $L$-Lipschitz in total variation of
the belief,
\begin{equation}
  \big|V^{\pi(h_t)}-V^{\pi(u_t)}\big| \;\le\; L\,\E\|b_t-\hat b_t\|_{TV}
  \;\le\; L\sqrt{\tfrac{1}{2}\varepsilon_t}
  \qquad\text{(Pinsker).}
  \label{eq:suff}
\end{equation}
The cross-entropy training objective directly minimizes an upper bound on
$\varepsilon_t$ at the readout (maximizing the $I(y;u_t)$ lower bound), so
end-to-end training \emph{co-adapts} the membrane to be the policy's
belief state rather than assuming it.
\end{proposition}
\begin{proof}
The identity $\E\,\mathrm{KL}(b_t\|\hat b_t)=I(y;h_t)-I(y;u_t)$ follows by
expanding both mutual informations against the common prior and using the
tower rule with $\hat b_t=\E[b_t\mid u_t]$; Pinsker converts KL to TV;
Lipschitzness transfers TV to values. DPI gives $\varepsilon_t\ge0$.
\end{proof}

\noindent\textbf{Empirical verification.} \texttt{V6\_sufficiency.py}:
(a) linear-probe gap between $u_t$ and the full membrane history estimates
$\varepsilon_t$ operationally; (b) with the backbone held fixed, the
membrane-driven policy beats random and fixed orders in anytime accuracy
--- the value of the information carried by $u_{t-1}$, isolated.
\texttt{A5\_policy} runs the trained-per-policy version at scale;
\texttt{S3\_forced\_start} shows closed-loop recovery after adversarial
forcing, which no open-loop (geometry-only) policy can express.

% -----------------------------------------------------------------------
\subsection{Why $d_{\ssp}$ barely matters and 2K parameters are enough}

\begin{theorem}[Bilinear rank cap]\label{thm:rank}
Eq.~\eqref{eq:score} depends on $(W_k,W_q)$ only through
$B=W_k^{\top}W_q\in\R^{D\times 6}$ with
$\operatorname{rank}(B)\le\min(d_{\ssp},6,D)$. Hence for every
$d_{\ssp}\ge 6$ the realizable score-function class is \emph{identical}: any
SSP with $d_{\ssp}=64$ is exactly reproduced by one with $d_{\ssp}=6$, and
the minimal exact parametrization has $6(D+6)$ parameters. A rank-$r$
factorization of $W_k$ with $r\ge 6$ is also exact; at $D=128$,
$d_{\ssp}=64$, $r=8$ this yields $8\cdot128+8\cdot64+64\cdot6=1{,}920\approx
2\mathrm{K}$ parameters, matching the paper's budget with zero loss.
\end{theorem}
\begin{proof}
$\mathrm{score}(m)=u^{\top}Bg_m/\sqrt{d_{\ssp}}$ and
$\operatorname{rank}(W_k^\top W_q)\le\min(\operatorname{rank}W_k,
\operatorname{rank}W_q)\le\min(d_{\ssp},6,D)$. Conversely any rank-$\le 6$
$B$ factors as $W_k^{\top}W_q$ with $d_{\ssp}=6$ (SVD), so the classes
coincide; parameter counts follow by arithmetic.
\end{proof}

\begin{corollary}[Prediction for the $d_{\ssp}$ ablation]\label{cor:dssp}
Task metrics must \emph{plateau} for $d_{\ssp}\ge 6$; any residual
differences are optimization-side (overparametrized linear factorizations
change the implicit-gradient conditioning; cf.\ Arora et al., 2018), not
representational. This converts the A4 ablation from a hyperparameter sweep
into a theorem check, and licenses shipping the 2K-parameter SSP.
\end{corollary}

\noindent\textbf{Empirical verification.} \texttt{V4\_rank\_equivalence.py}
verifies $\operatorname{rank}(B)\le6$ on the trained SSP and that the
rank-6 compression reproduces all scores (float tolerance) and $100\%$ of
argmax selections; \texttt{A4\_dssp} sweeps
$d_{\ssp}\in\{2,4,6,8,16,32,64,128\}$ plus the rank-8 budget variant and
logs parameters, accuracy, and selection agreement (Kendall $\tau$).

% -----------------------------------------------------------------------
\subsection{Why the $-\infty$ mask is not a detail}

\begin{theorem}[Coverage with masking; collapse without]\label{thm:mask}
With visited-anchor masking, the selection sequence samples without
replacement: after $t$ steps exactly $t$ distinct slices are observed
(full coverage at $t=K$). Without masking, consider fixed slice features
and the membrane recursion $u_{t+1}=\alpha u_t+(1-\alpha)x_{m_t}$ with
$\alpha\in(0,1)$. If a slice $m^\ast$ satisfies the \emph{absorbing
condition} $m^\ast=\arg\max_m \mathrm{score}(u^\ast,g_m)$ at its own fixed
point $u^\ast=x_{m^\ast}$, then once selected with membrane in the basin of
$u^\ast$, the policy reselects $m^\ast$ forever: the number of distinct
slices visited is bounded by the entry time, margins stall
($\E[\Delta M_t]\to0$), and both accuracy and expected exit latency degrade
simultaneously.
\end{theorem}
\begin{proof}
Without masking, repeated selection of $m^\ast$ makes
$u_t\to u^\ast=x_{m^\ast}$ geometrically at rate $\alpha$ (Banach
contraction). Scores are continuous in $u$, so if the argmax at $u^\ast$ is
strictly $m^\ast$, there is a neighborhood of $u^\ast$ where the argmax is
constant; the iteration enters it in finite time and never leaves
(absorbing). A constant observation adds no new evidence, so logits and
margins converge and increments vanish; with drift $\to0$,
Theorem~\ref{thm:stop} predicts diverging $\E[T_\theta]$ (censoring), while
missing regions cap attainable accuracy at that of the absorbed subset.
\end{proof}

\begin{remark}[The energy paradox]
Removing the mask is sometimes proposed to ``let the policy re-attend.''
Theorem~\ref{thm:mask} predicts the opposite of a tradeoff: \emph{both}
lower accuracy \emph{and} more slices to exit --- a rare double dissociation
that is cleanly falsifiable, and a strong rebuttal asset.
\end{remark}

\noindent\textbf{Empirical verification.} \texttt{V5\_masking\_dynamics.py}
(same trained weights, mask off at inference: revisit rate, coverage
plateau below $K$, drift collapse, joint accuracy/latency degradation);
\texttt{A2\_masking} retrains without the mask for the stronger end-to-end
claim on every benchmark.

% -----------------------------------------------------------------------
\subsection{Why adaptivity dominates the whole frontier, not one point}

\begin{theorem}[Pareto dominance from pathwise margin dominance]\label{thm:pareto}
Let two policies $\pi_1,\pi_2$ induce margin processes $M^{1},M^{2}$ on the
same inputs. If $M^{1}$ stochastically dominates $M^{2}$ at every step
($\Prob(M^1_t>x)\ge\Prob(M^2_t>x)\;\forall x,t$), then for \emph{every}
threshold $\theta$: $T^1_\theta\preceq_{st}T^2_\theta$ (earlier exits), and
with exit accuracy governed by Theorem~\ref{thm:risk} at equal $\theta$,
$\pi_1$'s (accuracy, expected-slices) curve weakly dominates $\pi_2$'s at
all operating points --- i.e.\ dominance of the entire Pareto frontier, not
a single-$\theta$ comparison.
\end{theorem}
\begin{proof}
By coupling, realize $M^1,M^2$ on a common probability space with
$M^1_t\ge M^2_t$ pathwise (Strassen's theorem, applied stepwise). First
passage over any level $\theta$ is antitone in the path, so
$T^1_\theta\le T^2_\theta$ pathwise, giving stochastic dominance of exit
times and $\E[T^1_\theta]\le\E[T^2_\theta]$; the risk bound at exit depends
only on $\theta$, so matched-$\theta$ accuracy floors coincide while cost
strictly improves wherever dominance is strict.
\end{proof}

\begin{remark}
The premise (stepwise margin dominance of SSP over fixed/random order) is
an \emph{empirical condition we test}, not an assumption smuggled into the
conclusion: one-sided Kolmogorov--Smirnov tests on per-step margin CDFs and
exit-time CDFs. Where dominance provably grows --- occlusion and non-uniform
density making region informativeness heterogeneous --- the ASP advantage is
predicted to \emph{widen} (tested by S1/S2).
\end{remark}

\noindent\textbf{Empirical verification.} \texttt{A5\_policy} (exit-time
CDFs + per-step margin CDFs, one-sided KS); \texttt{A1\_theta} traces the
frontier; \texttt{S1\_occlusion}/\texttt{S2\_density} test the widening
prediction under corruption severity sweeps on point clouds and CIFAR
patch-slices.

% -----------------------------------------------------------------------
\subsection{Training/inference consistency of Gumbel selection}

\begin{proposition}[Exact selection-consistency law]\label{prop:gumbel}
With i.i.d.\ Gumbel noise $g$, training-time hard sampling selects
$\arg\max_m(\mathrm{score}(m)+g_m)$, and by the Gumbel-max identity
\begin{equation}
  \Prob\big[\text{train-time selection}=\text{inference argmax}\big]
  \;=\;\mathrm{softmax}(\mathrm{score})_{\max}.
  \label{eq:gumbelmax}
\end{equation}
The temperature $\tau$ does not affect \eqref{eq:gumbelmax} (hard path);
it only scales the straight-through backward weights, whose bias vanishes
as $\tau\to0$ at the cost of gradient variance --- hence the annealing
schedule $1.0\to0.5$. As training sharpens score gaps,
$\E[\mathrm{softmax(score)}_{\max}]\to1$: the train/infer selection gap
closes \emph{automatically}, with no distillation or correction term.
\end{proposition}

\noindent\textbf{Empirical verification.} \texttt{V7\_gumbel.py} checks
identity~\eqref{eq:gumbelmax} to $\pm3$pp against Monte-Carlo Gumbel draws
on trained and untrained SSPs (trained $\gg$ untrained agreement), and
verifies finite nonzero straight-through gradients at
$\tau\in\{1.0,0.5,0.1\}$.

% -----------------------------------------------------------------------
\subsection{Geometry descriptors as information proxies}

\begin{proposition}[Spread is a max-entropy saliency prior]\label{prop:geom}
Model region $m$'s points as Gaussian with covariance $\Sigma_m$. The
region's differential entropy is
$h_m=\tfrac12\log\big((2\pi e)^3|\Sigma_m|\big)$, monotone in the
intra-cluster spread (for isotropic $\Sigma_m=\sigma_m^2 I$,
$h_m=3\log\sigma_m+\mathrm{const}$). At $t{=}0$, $u_0=0$ carries no class
evidence, so the optimal myopic choice reduces to a \emph{static}
max-entropy rule --- rank by spread; as $t$ grows, the posterior sharpens
and membrane-dependent (class-conditional) terms should dominate. Two
falsifiable predictions follow: (i) ablating \texttt{spread} damages the
\emph{first} selections most (largest early-drift loss in A3); (ii) the
gradient attribution of scores onto descriptor components shifts from
spread-dominant to centroid/dist-dominant over $t$. Under non-uniform
density (S2), the normalized-count component is the reliability weight that
suppresses sparse slices; dropping it should erase the S2 robustness gap.
\end{proposition}

\noindent\textbf{Empirical verification.} \texttt{A3\_geometry} (component
knockouts incl.\ \texttt{only\_spread}; per-step drift; selection Kendall
$\tau$ vs.\ full descriptor), \texttt{S2\_density} ($\pm$\texttt{count}
cross), per-timestep score-gradient attribution logged by the runner.

% -----------------------------------------------------------------------
\subsection{Claim--evidence map}
\begin{table}[h]
\centering\small
\begin{tabular}{@{}llll@{}}
\toprule
Result & Verified by & Scale experiment & Prediction falsified if \\
\midrule
Thm~\ref{thm:risk} & V2 & A1, S5 (all datasets) & risk$_{\mathrm{exited}} > \frac{(C-1)(1-\theta)}{C}$ \\
Thm~\ref{thm:stop} & V1 & A1, S6 & $\E[T_\theta]$ above \eqref{eq:wald}; no $\theta$-affinity \\
Cor~\ref{cor:bimodal} & V1 & S6 & unimodal exit histogram \\
Thm~\ref{thm:submod} & V3 & A5 anytime & rising marginal gains; high violation rate \\
Cor~\ref{cor:greedy} & V3 & A5 & SSP no closer to greedy than random \\
Prop~\ref{prop:suff} & V6 & A5, S3 & probe gap large; no recovery in S3 \\
Thm~\ref{thm:rank} & V4 & A4 & accuracy gains for $d_{\ssp}\!>\!6$ beyond noise \\
Thm~\ref{thm:mask} & V5 & A2 & no-mask improves accuracy or latency \\
Thm~\ref{thm:pareto} & A5 KS tests & A1, S1, S2 & fixed-order exits dominate anywhere \\
Prop~\ref{prop:gumbel} & V7 & --- & identity \eqref{eq:gumbelmax} violated \\
Prop~\ref{prop:geom} & A3 & S2 & spread knockout hurts late, not early \\
\bottomrule
\end{tabular}
\caption{Every theoretical claim is attached to an executable test and an
explicit falsification condition.}
\end{table}

\subsection{Scope and limitations}
Theorems~\ref{thm:risk} and \ref{thm:stop} condition on measurable
properties of the trained network (calibration; positive drift) that we
audit and report rather than assume silently; Theorem~\ref{thm:submod}
uses a conditional-independence surrogate whose violation rate is measured;
Theorem~\ref{thm:pareto}'s premise is itself an empirical test. We regard
this assumption-auditing loop --- theory states a measurable premise, the
suite measures it, the corollary then binds --- as the correct standard of
``evidence-backed theory'' for an applied venue.

\end{document}
